\subsection{Orthogonal ternary points and quadratic induction in degree four}
\label{hr:section}

Work with the marked representations in Propositions~\ref{spreal:strong},
\ref{spreal:four} and~\ref{dp:curvebranches}. The coefficient field \(E\)
below is a finite extension of \(\mathbb Q_\ell\) for an integral
arithmetic point, or the fraction field of a normalized test curve in
characteristic \(\ell\). Its residue field is finite. We allow finite
extensions of \(E\). Statements about classical attachment use
geometric Frobenius and the normalization
\(\operatorname{rec}(\pi_v|\det|^{(1-n)/2})\) in degree \(n\).

\begin{proposition}[The odd orthogonal ternary branch]\label{hr:orth}
Suppose that \(\rho\) on \(G_K\) is absolutely irreducible and is
isomorphic to its own dual. Let \(S\) be its invariant bilinear form
and \(J_3\) the antidiagonal unit matrix in the marked basis. The
projective class of \(J_3S\) has a unique representative \(C\) in
\(\mathrm{SO}(S)\). Suppose \(C\ne1\). In equal characteristic
such a point cannot occur. In characteristic zero there is a
regular algebraic, conjugate self-dual cuspidal representation
\(\Pi\) of \(\mathrm{GL}_3(\mathbb A_K)\) whose associated
representation is \(\rho\varepsilon^{-1}\).
\end{proposition}
\begin{proof}
The form is unique up to scalar, nondegenerate and symmetric by
irreducibility and odd dimension. The determinant of \(\rho\)
squares to one and reduces to one, so is one. The marked polarization
is \(\rho(chc^{-1})=J_3\rho(h)^{-t}J_3\). Consequently \(D=J_3S\)
implements conjugation by \(c\), \(D^2\) is scalar and
\(D^t S D=\mu S\). Since \((\det D)^2=\mu^3\), the required
representative is \(C=(\mu/\det D)D\), which is independent of the
scale of \(S\). Its square is a scalar in \(\mathrm{SO}_3\), hence
is one. This extends \(\rho\) continuously to \(G_{\mathbb Q}\)
in \(\mathrm{SO}(S)\) by sending \(c\) to \(C\).

Split \(S\) over a finite extension. Identify its special orthogonal
group with \(\mathrm{PGL}_2\) on trace-zero matrices, with quadratic
form their determinant. Indeed its isotropic conic consists of the
nilpotent lines indexed by their kernels in \(\mathbb P^1\).
Conjugation by \(\mathrm{PGL}_2\) induces every automorphism of this
conic, and a linear isometry fixing it pointwise is projectively
scalar, hence is the identity in \(\mathrm{SO}_3\). Write \(P\) for
the resulting projective representation.

We construct a linear lift over a DVR. After adjoining square roots
of representatives of \(E^\times/(E^\times)^2\), the compact image
of \(P\) lifts setwise to \(\mathrm{SL}_2\). The square-class
group is finite by Hensel's lemma and odd residue characteristic.
The inverse image of the compact group under this finite morphism
is a compact group: local sections are obtained by fixing one
nonzero entry of a matrix representative and taking a continuous
square root of its determinant near a given lift. A compact group
of matrices preserves a lattice, by taking the integral span of
the translates of one lattice. In particular we can reduce \(P\)
projectively on this lattice.

The normal subgroup \(H=\ker(\bar\Psi|_{G_K})\) has pro-\(\ell\)
image in the original ternary matrices, since all its matrices
reduce to the identity. Its projective reduction is a normal
\(\ell\)-subgroup. If nontrivial, it fixes a unique line over the
algebraic closure of the residue field (use the upper unitriangular
Sylow subgroup of \(\mathrm{PGL}_2\) over a finite field).
If trivial, the residual projective image is abelian, since
\(G_{\mathbb Q}/H\) embeds in the product of the images of \(\eta\)
and \(\bar\Psi\). Some inertia element has residual eigenvalue
ratio of order \(>2\): its adjoint characteristic polynomial on
the original lattice has roots \(1,z,z^{-1}\) with
\(z=\bar\Psi(s)\), and \(\bar\Psi^2|_{I_\ell}\ne1\).
Characteristic polynomials of the adjoint are independent of
lattice. The centralizer of this projective residual element is a
torus. In both cases the projective reduction is triangular after
a finite residue extension.

It can be made diagonal. In an upper triangular residual basis,
extend ramification by adjoining a square root \(u\) of the
uniformizer and conjugate by \(\operatorname{diag}(u,1)\).
The lower left entries were divisible by \(u^2\). Write
\(\bar\lambda\) for the residual diagonal ratio. The nontrivial
involution \(P(c)\) cannot be in the pro-\(\ell\) congruence
kernel, so \(\bar\lambda(c)=-1\). In equal characteristic take the
constant character \(\lambda=\bar\lambda\). In characteristic zero
take
\[
 \lambda=\widetilde{\bar\lambda}
                 (\varepsilon/\widetilde{\bar\varepsilon})^e,
\]
where the tilde denotes Teichm\"uller lift. For each \(g\) there is
exactly one representative \(W(g)\) with determinant \(\lambda(g)\)
and reduction \(\operatorname{diag}(\bar\lambda(g),1)\). Indeed
first arrange the specified reduction by scalar multiplication,
then correct the determinant by the unique square root in the
principal units. Uniqueness gives multiplicativity, and the same
construction on coordinate neighborhoods gives continuity. The lift
is odd and unramified outside finitely many places: where the
projectivization, determinant and residual characters are
unramified, its inertia is scalar, squares to one and reduces to
one. We have \(\operatorname{Ad}^0W|_{G_K}\simeq\rho\).

In equal characteristic the flag of Proposition~\ref{dp:curvebranches}
triangularizes inertia on \(\rho\). The inertia subgroup with all
diagonal entries one consists precisely of its unipotent elements.
It is therefore normal under the whole decomposition group, acting
through the extended ternary representation. Its image under \(P\)
consists of projective unipotents. If nontrivial they fix a unique
common line. Indeed two with distinct fixed lines have lifts
\(\left(\begin{smallmatrix}1&b\\0&1\end{smallmatrix}\right)\) and
\(\left(\begin{smallmatrix}1&0\\d&1\end{smallmatrix}\right)\)
with \(bd\ne0\) in the basis of those lines. The product of the
first and the \(j\)-th power of the second has trace \(2+jbd\),
which for some \(j\) differs from both \(2\) and \(-2\).
Normality then gives a decomposition-group line. If the subgroup
has trivial image, inertia is finite abelian by the finite
diagonal characters on the flag and faithfulness of the adjoint
action on \(\mathrm{PGL}_2\). It preserves the two eigenspaces of
the inertia element with residual ratio of order \(>2\). They
reduce to distinct lines in the diagonal residual basis, so the
decomposition group, which normalizes inertia and preserves each
residual line, cannot interchange them. The inertia eigenvalue
ratios have finite order at every place, as do the determinants,
so both inertia eigencharacters have finite image. Lemma~\ref{dp:pan}
now makes \(W\) geometrically reducible; its adjoint on \(G_K\)
would preserve the nilpotent line associated to an invariant line.

In characteristic zero the eigenline of the adjoint with inertia
character \(\varepsilon^e\) on \(I_\ell\cap\ker\psi_\ell\) is
unique by Lemma~\ref{spreal:flag}. It is preserved by the full
decomposition group since the character is conjugation invariant.
Its elements are nilpotent: their determinant is fixed by
conjugation and is multiplied by \(\varepsilon^{2e}\) under the
specified scaling. The kernel of any nonzero element therefore
gives a decomposition-group line of \(W\). The ratio of the
subcharacter to the quotient character on this open inertia
subgroup is \(\varepsilon^e\). After further finite restriction
the determinant there is \(\varepsilon^e\), so the ordered
characters there are \((\varepsilon^e,1)\). Both local characters
are de Rham, the upper one of strictly smaller Hodge--Tate weight
in the convention \(\mathrm{HT}(\varepsilon)=-1\). The residual
ratio is nontrivial on inertia. The distinguished ordinary
modularity theorem of Pan applies over \(\mathbb Q\) with the
specified odd determinant, to the maximal ideal defined by \(W\)
in the pseudocharacter ring with \(\ell\) inverted. It is maximal
because the integral image is a submodule of the finite coefficient
DVR over \(\mathbb Z_\ell\), and its localization is a field.
Irreducibility follows from that of the adjoint on \(G_K\).
If Pan writes \(W=W_f\otimes\theta\), its upper character and
the character of the same stable line in the de Rham representation
\(W_f\) are de Rham, so their ratio \(\theta_\ell\) is de Rham.
It is a cyclotomic integer power on open inertia. Away from
\(\ell\), the determinant identity shows that \(\theta^2\), hence
\(\theta\), has finite inertia. Dividing by the cyclotomic power
leaves a finite-order character by global class field theory.
Thus \(W\) is classical with its algebraic normalization.
\cite[Theorem~5.1.1(2)]{panFM}

Let \(\pi\) be its cuspidal representation of \(\mathrm{GL}_2\).
Its quadratic base change to \(K\) is cuspidal, since the
alternative gives a quadratic self-twist, hence reducibility of
\(W|_{G_K}\) by the two eigenspaces of the twist intertwiner.
The base change itself has no nontrivial self-twist: one would be
quadratic by the central character and split \(W\) on its open
kernel, contrary to strong irreducibility of the adjoint
(Proposition~\ref{spreal:strong}). These implications use equality
of unramified Frobenius polynomials and Chebotarev. Apply
Gelbart--Jacquet to the unitary normalization of this base change;
the adjoint parameter is unaffected by any common scalar twist.
The resulting cuspidal
adjoint lift is self-dual and invariant under
\(\operatorname{Gal}(K/\mathbb Q)\). Its archimedean parameter has
three distinct exponents \(e,0,-e\). At unramified places its
Satake parameters are the eigenvalues of \(\operatorname{Ad}^0W\):
the common scalar in the degree-two normalization cancels. The
degree-three arithmetic normalization multiplies them by
\(|\cdot|^{-1}\), giving \(\rho\varepsilon^{-1}\) as claimed.
\cite[Chapter~III, Theorems~4.2 and~5.1]{arthurClozelBC}
\cite[Theorem~9.3]{gelbartJacquetAdj}
\end{proof}

\begin{corollary}[Essential self-duality]\label{hr:ess}
For a strongly irreducible ternary point, an isomorphism
\(\rho^\vee\simeq\rho\otimes\gamma\) on \(G_K\) already implies
\(\gamma=1\). If its orthogonal extension has \(C=1\), it too has
a linear adjoint lift over \(\mathbb Q\), this time with scalar
complex conjugation, and with the same ordinary local description
as in Proposition~\ref{hr:orth}.
\end{corollary}
\begin{proof}
In characteristic zero the determinant of \(\rho\) is one by the
determinant calculation in Proposition~\ref{spreal:end}, which
uses the local flag, the residual inertia kernels and
\(\ell\nmid h_K\) without any use of the block decomposition.
In equal characteristic it follows directly from the finite inertia
character argument and its trivial reduction. Taking determinants
therefore gives \(\gamma^3=1\). Applying \(c\) to
\(\rho^\vee\simeq\rho\otimes\gamma\) and using
\(\rho^c\simeq\rho^\vee\) shows that \(\gamma\gamma^c\)
is a self-twist. Every finite-order self-twist of a strongly
irreducible representation is trivial: its intertwiner is scalar
on the open kernel by Schur's lemma. Hence \(\gamma^c=\gamma^{-1}\).
The residual characters of \(\rho\) are
\(1,\bar\Psi,\bar\Psi^{-1}\); comparison of the two residual
multisets forces \(\bar\gamma\) to be one of them. All three are
fixed by \(c\). Thus \(\bar\gamma^2=1\) and \(\gamma=1\), since
reduction is injective on the cube roots for \(\ell\ge5\).
If \(C=1\), the construction of \(P,W,\lambda\) and their local
flags is unchanged, with \(\bar\lambda(c)=1\) and \(W(c)=1\).
This reduces that case to an even two-dimensional problem;
the odd-determinant theorem of Pan used above does not settle it.
\end{proof}

\begin{proposition}[A real quadratic self-twist in degree four]\label{hr:ind}
Let \(R\) be an absolutely irreducible point of the good rank-four
problem admitting a nontrivial even quadratic self-twist \(\delta\).
Put \(a=\bar\alpha=\bar\xi\) and \(b=\bar\Psi\). Then \(b=\bar\delta\)
and the inducing field \(M\) is the real quadratic field cut out by
\(b\). In equal characteristic such a point cannot occur if
\(\ell\) is unramified in \(M\). In characteristic zero \(\ell\)
necessarily splits in \(M\), and \(R\) has a regular cuspidal
attachment on \(\mathrm{GL}_4/\mathbb Q\). Its base change to any
imaginary quadratic field on which \(R\) remains irreducible is
cuspidal and has the restricted Frobenius parameters.
\end{proposition}
\begin{proof}
A self-twist intertwiner has scalar square. After scaling it to
square one, its two eigenspaces are preserved by \(G_M\) and
interchanged by the other coset, so both have dimension two. Thus
\(R=\operatorname{Ind}_{G_M}^{G_{\mathbb Q}}U\), with \(U\)
irreducible. Conversely any such quadratic induction admits this
self-twist. On residual semisimplifications multiplication by
\(\bar\delta\) preserves the multiset \(\{1,a,b,ab\}\).
The only possibilities for the image of the trivial character are
the even entries \(1,b\); the quadratic character survives
reduction since \(\ell\) is odd. Over \(G_M\) the residual multiset
is consequently \(1,1,a,a\). The two conjugate blocks have
identical residual semisimplifications: their characters are
chosen from \(1,a|_{G_M}\), both fixed by conjugation from
\(G_{\mathbb Q}\). This follows by reducing lattices in their
direct sum and using independence of the residual semisimplification.
Thus each block has residual characters \(1,a|_{G_M}\).
They have opposite signs at both real places of \(M\).

Intersections of the ordinary local flag with the plane of \(U\)
give a full flag after omitting repeated subspaces. This applies
at either place of \(M\) over \(\ell\) by using both conjugate
planes. In equal characteristic the inertial eigencharacters are
finite at every finite place. The determinant of \(U\) is therefore
the constant character \(a|_{G_M}\), by the character argument
preceding Lemma~\ref{dp:pan}. Apply the proof of that lemma over
\(M\). This is an abelian totally real field unramified at
\(\ell\), the residual ratio extends to \(G_{\mathbb Q}\), and
the set of places may contain every ramified place and those over
\(\ell\). At the places over \(\ell\) the residual ratio is either
nontrivial at all of them or trivial at all, since it is the
restriction of a character of \(G_{\mathbb Q}\) and \(M\) is
Galois. In the latter case use the actual upper character at each
place in the character-labelled ring. The map from the ordinary
weight algebra kills \(\ell\) and each local pro-\(\ell\) inertia
augmentation, including when a place has residue degree two.
Finiteness over this algebra therefore makes every
characteristic-polynomial coefficient constant. Brauer--Nesbitt
contradicts irreducibility of \(U\).
\cite[Theorems~5.1.1(1) and~6.1.1(1)]{panFM}

In characteristic zero the four distinct exponents of \(R\) on
an open subgroup of inertia are \(e+1,e,1,0\), in descending
order. On a common open subgroup of \(I_\ell\cap G_M\) the two
planes have disjoint pairs of exponents. A decomposition element
interchanging them would preserve these exponents, since
\(\varepsilon\) is a global character. Thus \(\ell\) splits in
\(M\). The exponents of \(\det U\) at its two places above \(\ell\)
are the sums of the two respective pairs. Global reciprocity on a
power of a nontrivial totally positive norm-one unit of \(M\)
equates these sums, by the calculation in Lemma~\ref{spreal:even}.
Their common value is \(e+1\), and the pairs are
\(\{0,e+1\},\{1,e\}\). The character \(\det U/\varepsilon^{e+1}\)
has finite inertia at all places and is unramified outside finitely
many. A common power is everywhere unramified and factors through
the finite narrow class group; hence this character has finite
image.

The local upper character of \(U\) has the larger cyclotomic
exponent in each pair and hence strictly smaller Hodge--Tate
weight than the quotient. Both characters are de Rham, since
they have the indicated powers on open inertia. The residual
ratio \(a\) is distinguished there: the chosen Frobenius lift
\(\sigma_\ell\) has \(a(\sigma_\ell)=\bar\xi(\sigma_\ell)\ne1\)
by the good-datum hypothesis and the splitting of \(\ell\).
The ordinary theorem now makes \(U\) Hilbert modular over \(M\),
up to the standard algebraic twist. The coefficient map gives a
maximal ideal after inverting \(\ell\), as in
Proposition~\ref{hr:orth}. To absorb the twist in \(U=W_f\otimes\theta\),
compare the upper character of \(U\) at each place above \(\ell\)
with that of the same line in the de Rham \(W_f\). Their ratio
\(\theta_v\) is de Rham and hence a cyclotomic integer power on
open inertia at the split place. The determinant identity gives
finite inertia for \(\theta\) at all other finite places.
The same norm-one unit argument equates the two exponents; dividing
by their common cyclotomic power leaves a finite-order character.
Both twists thus have algebraic automorphic counterparts.
\cite[Theorem~5.1.1(2)]{panFM}

Write \(\pi_M\) for the corresponding cuspidal representation in
the normalization giving \(U\). Its Galois conjugate has the other
pair of exponents at the first place above \(\ell\), so the two
are inequivalent. Cyclic automorphic induction gives a cuspidal
representation, and the arithmetic normalization in degree four
is obtained by taking
\(\Pi=\operatorname{AI}_{M/\mathbb Q}(\pi_M)|\det|\).
Indeed \(1-3/2=-1/2\), the shift used for \(U\) in degree two.
The unramified parameters therefore give \(R\). At the real place
the parameter is the sum of those at the two real places of
\(M\), with the same shift of normalization; its four arithmetic
exponents are \(0,1,e,e+1\), giving regular algebraicity.
If quadratic base change to \(K\) failed to be cuspidal,
\(\Pi\) would admit its quadratic self-twist. The unramified
polynomials would give the same self-twist of \(R\) and split
\(R|_{G_K}\) by the intertwiner argument above.
\cite[Chapter~III, Theorems~4.2, 5.1 and~6.2, Corollary~6.5]{arthurClozelBC}
\end{proof}

\begin{corollary}[Imprimitive quartic points]\label{hr:imprimitive}
Every absolutely irreducible characteristic-zero point of the good
quartic problem which fails strong irreducibility is induced from a
plane over the real quadratic field cut out by \(b=\bar\Psi\).
It has the cuspidal attachment of Proposition~\ref{hr:ind} and
its finite symplectic adjoint Selmer group over \(\mathbb Q\)
vanishes. No such point permutes a system of four lines.
\end{corollary}
\begin{proof}
First consider any self-twist \(R\simeq R\otimes\delta\).
Determinants give \(\delta^4=1\). Over the decomposition group
at \(\ell\) the flag of Proposition~\ref{spreal:four} has
characters \(\kappa_i\). On an open subgroup of inertia in
\(\ker\delta\) they are four distinct cyclotomic powers.
Thus the permutation matching \(\{\kappa_i\delta\}\) with
\(\{\kappa_i\}\) fixes each character, and
\(\delta|_{G_{\mathbb Q_\ell}}=1\).
The image of the trivial residual constituent under twisting gives
\(\bar\delta\in\{1,a,b,ab\}\). The specified element \(s_\ell\)
of this decomposition group satisfies
\(a(s_\ell)b(s_\ell)^i\ne1\) for all integers \(i\). This
eliminates \(a,ab\). Reduction is injective on the fourth roots of
unity, so the self-twist group has order at most two, and its sole
possible nonidentity element is even quadratic with reduction \(b\).
\src{src/Development/AutomorphicArithmetic/AutomorphicVarieties/SymplecticGSp4/DecompSigmaOrd.lean}{114}{The local character condition of the good datum}.

For four permuted lines choose the direct sum of their intersections
with a stable lattice. The nonzero entries of the resulting monomial
matrices are units, so their permutation quotient \(P\) survives
reduction. Independence of residual semisimplification gives an
invariant residual flag with one-dimensional quotients. The derived
group of the finite residual image is therefore an \(\ell\)-group,
as it is contained in the upper unitriangular matrices. Hence
\(P'\) is a normal \(\ell\)-group in the transitive group \(P\).
Its orbits on the four lines all have the same size, a divisor
of four and a power of \(\ell\). Thus it fixes every line and is
trivial as a permutation group. This holds also for \(\ell=3\).
The faithful transitive abelian group \(P\) is regular of order
four. Index the lines by \(P\). Multiplying the line indexed by
\(g\) by \(\chi(g)\) gives a self-twist for each character
\(\chi\) of \(P\). These four distinct twists contradict the
bound of two.

Now restrict to an open normal subgroup
\(H\) on which the representation is reducible. The translates
of any simple \(H\)-subspace span \(R\), so the restriction is
semisimple and its isotypic subspaces are transitively permuted.
Every simple constituent has multiplicity one, because some
element of \(I_\ell\cap H\) has four distinct eigenvalues by the
four distinct cyclotomic exponents on open inertia. Thus the
remaining alternative is two inequivalent planes. Their permutation
gives the quadratic self-twist just identified, so
Proposition~\ref{hr:ind} applies. Choose an imaginary quadratic
\(K\) unramified at \(\ell\) with irreducible restriction by
Lemma~\ref{rk4:quadratic}. Its cyclotomic-kernel irreducibility
follows from Lemma~\ref{spreal:twists}, since the ordinary characters
remain distinct on open inertia and inertia surjects onto the
cyclotomic quotient. The cuspidal base change supplied by
Proposition~\ref{hr:ind} and the given similitude pairing now
satisfy Proposition~\ref{newgal:adselmer}(2). Its symplectic
direct summand has zero finite Selmer group. In particular every
tangent of Proposition~\ref{spreal:four} is a coboundary there.
\end{proof}

\begin{proposition}[Strong irreducibility on the quartic test curves]
\label{hr:case3strong}
For the original good quartic datum in its third host character case,
every absolutely irreducible point in either characteristic remains
absolutely irreducible over every finite number-field extension.
In the second host case an absolutely irreducible equal-characteristic
curve failing this condition requires \(b\) to cut out a real
quadratic field ramified at \(\ell\).
\end{proposition}
\begin{proof}
Write \(N\) for the host level and \(a=\bar\xi\),
\(b=\bar\varepsilon^e\bar\psi\). For a self-twist, determinants
give \(\delta^4=1\) in either characteristic. Its reduction is one
of \(1,a,b,ab\). We check the orders of these characters on the
original group.
\src{src/Development/AutomorphicArithmetic/AutomorphicVarieties/AutomorphicCases/SymplecticGSp4.lean}{15}{The conductor factors of the auxiliary character}
and \src{src/Development/AutomorphicArithmetic/AutomorphicVarieties/HeightAndVanishing/DirichletPrimeToLPartMulLPowerPart.lean}{221}{The exponents in the two host cases}.

In case three, \(m=\operatorname{ord}(\psi)>1\) is odd and prime to
\(\ell\). The auxiliary odd prime \(p_1\) divides neither \(N\)
nor \(\ell\), and \(\ell\nmid N\). The character \(\xi\) is the
product of a character of two-power conductor and a character of
order \(m\) modulo \(p_1\). Thus on inertia at \(p_1\) the
characters \(a,ab\) both have order \(m\): inertia surjects
onto \(\operatorname{Gal}(\mathbb Q(\mu_{p_1})/\mathbb Q)\),
and reduction preserves its prime-to-\(\ell\) character values.
Also the image of \(b\) on \(\ker\bar\varepsilon\) contains the
image of \(\bar\psi\), by the Chinese remainder theorem at
\(N,\ell\). None of these three characters has fourth power one.
\src{src/Development/AutomorphicArithmetic/AutomorphicVarieties/HeightAndVanishing/BchiFacts.lean}{8}{The coprime conductors in case three}
and \src{src/Development/AutomorphicArithmetic/AutomorphicVarieties/WeakShimuraTheory/DirichletGroupPowers/DirGPowOne.lean}{77}{The order prime to the coefficient prime}.

In case two \(\xi\) is primitive of conductor \(2^n\). Put
\(u=v_2(N)\ge2\). By construction \(2^n\ge\ell N\);
the odd part of \(N\) contains an odd host prime, so \(n\ge u+3\).
The subgroup \(1+2^u\mathbb Z_2\) modulo \(2^n\) is cyclic
of order \(2^{n-u}\). Its order-two subgroup is the kernel of
reduction modulo \(2^{n-1}\), where primitivity makes \(\xi\)
nontrivial. Thus \(a\) is faithful on that cyclic subgroup.
On the corresponding inertia at two, \(b=1\), since \(\psi\)
has level \(N\) and \(\ell\) is odd. Consequently neither
\(a\) nor \(ab\) has fourth power one. Reduction is injective
on the fourth roots of unity, so the self-twist group has size at
most two, and any nontrivial twist is the even quadratic lift of
\(b\). In case three it is trivial. Corollary~\ref{hr:imprimitive}
now gives the characteristic-zero assertion.

In equal characteristic let \(H\) be an open normal subgroup of
\(G_{\mathbb Q}\) on which the representation becomes reducible
over an algebraic closure of the coefficient field. Its restriction
is semisimple by the sum-of-translates argument above. Thus the
commutant of \(H\) has dimension greater than one, and its trace-zero
part \(W\) is nonzero. Conjugation on \(W\) has finite image since
\(H\) acts trivially. In any stable lattice the residual
semisimplification is a sum of characters. Triangularize its
residual action. The preimage of the upper unitriangular subgroup
in the integral image is pro-\(\ell\), by the successive congruence
quotients, and its quotient is diagonal and abelian. Hence
the finite group acting on \(W\) has a normal \(\ell\)-subgroup
with abelian quotient. A finite \(\ell\)-group acting in characteristic
\(\ell\) has nonzero invariants: use a central element of order
\(\ell\), take its fixed space, and induct on the quotient.
The abelian quotient has a common eigenline in that invariant space.
A nonzero endomorphism on this line intertwines the original
representation with a character twist and is invertible by
irreducibility. The character is nontrivial, since a scalar
endomorphism of trace zero is zero for odd \(\ell\) in dimension
four. This excludes the curve in case three. In case two it forces
the quadratic self-twist of \(b\); Proposition~\ref{hr:ind}
excludes it whenever \(\ell\) is unramified in the inducing field.
\end{proof}

There is a matrix criterion for this induction argument. Suppose
the absolutely irreducible \(R\), already symplectic with multiplier
\(\mu\), also preserves a symmetric form up to a character \(\tau\).
Comparing the two identifications with the dual gives a self-twist
\(\delta=\tau/\mu\). It is nontrivial by Schur's lemma since one
form is symmetric and the other alternating. For the two original
host cases the character calculation in
Proposition~\ref{hr:case3strong} forces case two and the real
quadratic character lifting \(b\). Thus the characteristic-zero
point is covered by Proposition~\ref{hr:ind}, and an
equal-characteristic point requires ramification of \(\ell\) in
that quadratic field.
In particular a genuine tensor product \(R=X\otimes Y\) of two
degree-two representations has a symmetric form with
\(\tau=(\det X)(\det Y)\), by tensoring their alternating determinant
pairings. Here \(\tau^2=\det R=\mu^2\), so the self-twist is
quadratic. These criteria concern the representation on the
original group; the existence of an auxiliary tensor-factor
tower alone supplies no such form there.

The automorphy assertions here concern classical cuspidal objects
and their unramified arithmetic parameters. The rank-four source
predicate also asks for a witness in its specified automorphic
model for every qualifying quadratic CM restriction.
\src{src/Development/AutomorphicArithmetic/AutomorphicVarieties/ModularityArgument/AutomorphyLifting/IsAutCsdA.lean}{16}{The two clauses of the quadratic-restriction predicate}.

\begin{lemma}[An unramified test on the residual field]\label{hr:class}
In either rank let \(F_0\) be the field of the kernel of the actual
diagonal residual representation (on \(G_K\) in rank three).
Suppose the \(\ell\)-part of the narrow class group of \(F_0\) is
cyclic. Every equal-characteristic test curve unramified at all finite
places on \(G_{F_0}\) has constant characteristic polynomials.
It suffices that the curve be tame above \(\ell\) over \(F_0\)
and have inertia image of order prime to \(\ell\) at the remaining
finite places over the allowed ramified primes.
\end{lemma}
\begin{proof}
The image of \(G_{F_0}\) lies in the first congruence group and is
pro-\(\ell\). Tame inertia alone can still have an \(\ell\)-part
away from \(\ell\); at those places the defining inertia kernels
must also kill that part. In rank three they kill the kernel of the
residual action. In rank four, at a prime away from \(\ell\) where
the prescribed condition only kills the kernel of \(\psi\), require
in addition that the order of its allowed inertia quotient be
prime to \(\ell\), or directly that the image on \(G_{F_0}\) be
trivial. At places over \(\ell\) tame inertia has pro-prime-to-\(\ell\)
order. Under these local conditions the pro-\(\ell\) image is
everywhere unramified. Any finite quotient is an \(\ell\)-group
whose abelianization is a quotient of the cyclic \(\ell\)-class
group, so is cyclic by the Frattini criterion. Its order is bounded
by the same class number. Passing to the inverse limit gives
finite image, also over the original number field since \(F_0\)
has finite degree. Every eigenvalue of every matrix has finite
order and is thus constant over the algebraic closure of the finite
residue field. Reduction recovers each characteristic polynomial.
\end{proof}

\subsection{Which ideal the tower controls}\label{hr:tower}

Write \(F\) for the real field of a fit, deep rank-four tower and
\(L/F\) for its quadratic CM extension. Let \(X.A\) be its framed
ring and \(X.B\hookrightarrow X.A\) the coefficient ring of the
characteristic polynomials and the ordered local characters. The
following finite-fibre calculation isolates the use of an ordinary
Hecke algebra in the source.

\begin{lemma}\label{hr:towerfin}
Suppose an auxiliary link supplies an ordinary determinant ring
\(D_1\) over its weight algebra \(\Lambda_1\), an ideal \(J_1\)
such that \(D_1/J_1\) is finite over \(\Lambda_1\), and the link
map \(f_1:D_1\to X.B\). This map identifies the determinant on
\(G_L\) and its local characters with the restricted ones twisted
by the link character \(\kappa\). Put \(O=c'.O\) and
\[
 U_{\mathrm{in}}=(\psi_{v,i}(s)-1\mid v\mid\ell,\ 
                   1\le i\le4,\ s\in I_v)\subset X.B .
\]
Then for some \(m\),
\[
 \mathfrak m_{X.B}^{m}
       \subset f_1(J_1)X.B+\mathfrak m_O X.B+U_{\mathrm{in}}.
\]
\end{lemma}
\begin{proof}
Pass to the quotient by any prime containing the right hand side.
It is a local domain of characteristic \(\ell\) with the fixed
finite residue field. The inertia weight generators of \(\Lambda_1\)
map to constants: after the fixed link twist the remaining ordered
inertia characters are one. The local units of \(O\), including
the values of the link twist, map to their residues. By the
universal property of the weight algebra its map is the residual
one, since that map has the same labelled characters. Finiteness
therefore makes the image of \(D_1/J_1\) integral over the finite
residue field. In a local domain with residue field \(k\), any
element integral over the coefficient field \(k\) belongs to \(k\):
the subalgebra it generates is a field and reduction embeds that
field into \(k\) over \(k\).

Hence the characteristic coefficients on \(G_L\) are constant
after removal of the scalar twist. For \(g\in G_F\), use
\(g^2\in G_L\). Each elementary symmetric polynomial in four
variables \(x_i\) is integral over the symmetric polynomials in
\(x_i^2\), since each \(x_i\) satisfies
\(\prod_j(T^2-x_j^2)=0\). The resulting monic identities with
integral symmetric coefficients hold on every matrix by the
characteristic polynomial, hence in \(X.B\) by its injection
into the framed ring. They survive the prime quotient without
any lifting of its prime to \(X.A\). Thus the characteristic
coefficients for \(g\) are constant there. Each marked decomposition
group lifts to the one on \(L\) by the splitting at these places;
the corresponding labelled characters, after division by the link
twist, are constant as well. These data topologically generate
\(X.B\), so the prime is its maximal ideal. In the noetherian ring
\(X.B\), the radical inclusion gives the asserted power.
\src{src/Development/AutomorphicArithmetic/AutomorphicVarieties/ModularityArgument/ThresholdArguments/DenseThresholdCases/TPmFin.lean}{126}{The prime-by-prime calculation for one link}.
\end{proof}

The ideal actually used is
\[
 J=f_1(J_1)X.B+f_2(J_2)X.B .
\]
The two links have the same twisting character \(\kappa\) and
distinct auxiliary primes. Lemma~\ref{hr:towerfin} needs only
the first link. The two links give separate local vanishing
statements at a seed and at the prescribed curve points. The
characteristic-zero automorphy assertions away from the separate
auxiliary primes are combined using their distinctness and the
common twist. The actual quantifier chooses a lower coefficient
extension for the tower, then provides such an ideal for every
larger coefficient DVR and every representing space \(X\).
\src{src/Development/AutomorphicArithmetic/AutomorphicVarieties/ModularityArgument/ThresholdArguments/DenseThresholdCases/LinkTwo.lean}{13}{Two links with a common twist};
\src{src/Development/AutomorphicArithmetic/AutomorphicVarieties/ModularityArgument/ThresholdArguments/DenseThresholdCases/PmTwo.lean}{21}{Sum of the two extended ideals}.

On the unitary side \(J_i\) is the unit ideal when the specified
ordinary module vanishes. Otherwise it is the kernel of a
surjection of the universal ordinary determinant ring onto the
localized Hecke algebra: the interpolated determinant gives the
map, the unramified Hecke polynomials and the ordinary uniformizer
eigencharacters recover the Hecke generators, and the map agrees
on the weight algebra. Finiteness of this Hecke algebra over the
weight algebra then proves finiteness of the quotient. Construction
of the interpolated determinant uses separation by arithmetic
points with ordinary Galois attachments. The additional hypotheses
in the modular ideal include density, the point theorem,
purity and classicality at the indicated points, twist compatibility
and the two component assertions.
\src{src/Development/AutomorphicArithmetic/AutomorphicVarieties/ModularityArgument/ThresholdArguments/DenseThresholdCases/ExistsHeckeIdeal.lean}{28}{The zero-module case and the Hecke kernel};
\src{src/Development/AutomorphicArithmetic/AutomorphicVarieties/ModularityArgument/ThresholdArguments/DenseThresholdCases/TJRing.lean}{172}{Transfer of ordinary Hecke finiteness};
\src{src/Development/AutomorphicArithmetic/AutomorphicVarieties/ModularityArgument/ThresholdArguments/DenseThresholdCases/THidaU.lean}{38}{Assembly of the unitary ideal properties}.

For a given curve the required extra assertion is that after
the specified finite base changes and an integral choice of
lattice there is \(y:X.A\to\mathcal O_{\rm curve}\) cutting out
its restricted representation, killing \(J\), and giving the
original labelled inertia characters. On a fit tower those
characters reduce to one on the marked inertia subgroups. The
lemma kills \(\mathfrak m_{X.B}\) under \(y\) since the target is
a domain. All characteristic polynomials on \(G_F\) are then the
residual ones. For every \(g\in G_{\mathbb Q}\) some positive
power belongs to \(G_F\). Every eigenvalue of \(g\) is thus a
root of unity, hence constant in equal characteristic; reduction
again gives the specified polynomial. This excludes every such
curve with a higher-dimensional simple constituent by
Brauer--Nesbitt. On the marked model the projectors used in
Proposition~\ref{dp:curvebranches} then give constancy of the
weight-zero entry monomials themselves.
\src{src/Development/AutomorphicArithmetic/AutomorphicVarieties/ModularityArgument/ThresholdArguments/DenseThresholdCases/PotentialModularityTower.lean}{39}{The finite-fibre, seed and component requirements};
\src{src/Development/AutomorphicArithmetic/AutomorphicVarieties/ModularityArgument/ThresholdArguments/DenseThresholdCases/ContradictionStarPromodular.lean}{135}{Recovering the residual characteristic polynomial over the tower}.

In the rank-four case-three construction, local vanishing at a
characteristic-zero seed gives inclusion on some component at its
residual point. A chain of linked residual frames carries this
to the residual point of the proposed test. At each good residual
point a separate component theorem promotes inclusion on some
component to inclusion on all of them. Only the latter conclusion
gives \(y(J)=0\) for the chosen lift. The source chooses a layer
with threshold \(30(C+1)^2\), where \(C\) is supplied by its
cohomology bound, and requires a seed over that layer and a
descending chain for the specific residual lattice. Its seed
local-vanishing proof also uses adjoint lifting. Merely producing
the ideal, or one characteristic-zero point of its zero set,
does not yield this quantifier over curves.
\src{src/Development/AutomorphicArithmetic/AutomorphicVarieties/ModularityArgument/ThresholdArguments/DenseThresholdCases/PromodularityCaseThreeOfLemmas.lean}{131}{The choice of layer, seed and chain};
\src{src/Development/AutomorphicArithmetic/AutomorphicVarieties/ModularityArgument/ThresholdArguments/DenseThresholdCases/BirthPm.lean}{28}{Adjoint lifting at a seed}.

In the rank-three construction there are separate ordinary
determinant ideals in degrees two and three. For each degree
the input supplies realizations with an auxiliary prime avoiding
any prescribed nonzero integer, and a seed specializing into
each ideal satisfying the point theorem. Component propagation
first treats degree two; degree three also uses the splitting
condition from the third component property and the degree-two
conclusion. It produces the degree-three finite-point assertion.
A second choice of auxiliary prime in each degree removes the
exceptional place from the two automorphy assertions.
\src{src/Development/AutomorphicArithmetic/AutomorphicVarieties/ModularityArgument/ThresholdArguments/DenseThresholdCases/PrLinkedEnd.lean}{22}{Specialization and the two modular ideals};
\src{src/Development/AutomorphicArithmetic/AutomorphicVarieties/CurveDVRSpecialization/AbsMain.lean}{12}{Propagation in the two degrees};
\src{src/Development/AutomorphicArithmetic/AutomorphicVarieties/ModularityArgument/ThresholdArguments/DenseThresholdCases/TwoRuns.lean}{80}{The two realizations and change of auxiliary prime}.

\paragraph{The ternary seed on the linked CM field.}
Here \(L\) is the CM field already chosen for the test. Write
\(\varepsilon\) for the cyclotomic character,
\(\Phi=\varepsilon^j\chi'\) for the prescribed lift of the residual
character and \(\eta\) for the auxiliary quadratic character. The
seed theorem requires \(\chi'\ne\eta\), \(M\ne0,\ 2\ell\mid M\),
the extension \(L/L^+\) unramified at every finite place and split
at all places dividing \(M\),
and triviality of \(\chi'|_{G_L}\) on inertia away from \(M\).
The linked field and character conditions include \(K\subseteq L\)
for the imaginary quadratic field of \(\eta\), \(L/\mathbb Q\) abelian,
and triviality of the auxiliary character \(\theta\) on \(G_L\).
In the seed datum the positive auxiliary conductor \(N'\) is prime to
the odd prime \(\ell\). The prime \(w\nmid N'\) acts trivially through
both \(\eta\) and \(\chi'\) on \(\mathbb Q(\mu_{N'})\), through which
they factor. The character \(\theta\) has order \(\ell\) and factors
through \(\mathbb Q(\mu_w)\). Also \(\chi'(c)=(-1)^j\) for complex
conjugation and \(\bar\Phi^2|_{G_L}\ne1\). The link and the congruence with
the prescribed level character give the two local hypotheses at
the point of use. For every lower bound \(b\) the theorem produces
\(m\ge b,\ m\equiv j\bmod(\ell-1)\), and a characteristic-zero
representation of \(G_L\) irreducible on open subgroups, with residual
characters \((\bar\Phi,1,\bar\Phi^{-1})|_{G_L}\), away-from-\(\ell\)
inertial types \((\chi',1,(\chi')^{-1})|_{G_L}\), ordered ordinary
exponents \((m,0,-m)\), a conjugate-dual pairing and automorphic
attachment after twisting by \(\varepsilon^{-1}|_{G_L}\).
Here one takes \(b=2\).
\src{src/Development/AutomorphicArithmetic/AutomorphicVarieties/IdeleGenerators/SlopeTestRemainingChoices.lean}{136}{Field and character conditions},
\src{src/Development/AutomorphicArithmetic/AutomorphicVarieties/ModularityArgument/ThresholdArguments/DenseThresholdCases/SymUse.lean}{27}{Ternary seed theorem}
and \src{src/Development/AutomorphicArithmetic/AutomorphicVarieties/ModularityArgument/ThresholdArguments/DenseThresholdCases/LinkedEndE.lean}{51}{Its use on the linked field}.

In the construction this is \(\operatorname{Ad}^{0}(\rho)|_{G_L}\)
for an auxiliary degree-two Eisenstein congruence. An upper triangular
action on a plane with diagonal \(a,b'\) induces on its traceless
endomorphisms a flag with diagonal \(a/b',1,b'/a\): take first the
strictly upper triangular endomorphisms, then all upper triangular
traceless endomorphisms. This proves the claimed exponents from
the difference \(m\) of the two ordinary exponents. The auxiliary
\(\eta,\theta\) disappear on \(G_L\). The trace pairing on traceless
endomorphisms is symmetric and invariant in characteristic zero.
Open irreducibility, existence of \(\rho\) and automorphic attachment
of its twisted adjoint are arithmetic inputs to this calculation.
\src{src/Development/AutomorphicArithmetic/AutomorphicVarieties/ModularityArgument/ThresholdArguments/DenseThresholdCases/AdjointLiftFromEisensteinSeedOfHyps.lean}{184}{Inputs and construction}
and \src{src/Development/AutomorphicArithmetic/AutomorphicVarieties/ModularityArgument/ThresholdArguments/DenseThresholdCases/AdjointLiftFromEisensteinSeedOfHyps.lean}{250}{Restriction of the adjoint and its local conditions}.

For either degree \(n=2,3\) the seed specialization is taken after
the link twist \(\kappa\). The integral comparison characters are
\(\Psi^{w_i}|_{G_L}\kappa\), with \(w_i=(1,-1)\) or \((1,0,-1)\)
and \(\Psi\) the original integral character. Their reductions are
the labelled residual characters in the determinant problem.
Away from \(\ell\) an inertial eigenvalue in the required kernel is
a root of unity \(u\) of order dividing a number \(h\) prime to \(\ell\),
with \(u\equiv1\) in the residue field. Since
\(1+u+\cdots+u^{h-1}\equiv h\ne0\), the identity \(u^h=1\)
gives \(u=1\). Diagonality on that inertia subgroup kills the
whole action there. Above \(\ell\) scalar twisting multiplies each
graded character by \(\kappa\) in the same flag, preserving their
order. Thus the representation and its labelled characters give a
point of the ordinary determinant ring. The cited arithmetic
point theorem puts this particular point in \(V(J)\).
\src{src/Development/AutomorphicArithmetic/AutomorphicVarieties/ModularityArgument/ThresholdArguments/DenseThresholdCases/PrSeedPtOf.lean}{132}{Integral comparison characters},
\src{src/Development/AutomorphicArithmetic/AutomorphicVarieties/ModularityArgument/ThresholdArguments/DenseThresholdCases/PrSeedPtOf.lean}{198}{Inertia and flags}
and \src{src/Development/AutomorphicArithmetic/AutomorphicVarieties/ModularityArgument/ThresholdArguments/DenseThresholdCases/PrSeedPtOf.lean}{219}{Application of the point theorem}.

Write \(B\) for this determinant ring and choose a minimal prime \(Z\)
of \(J\) below the kernel of the seed point. Put \(C=B/Z\). If the
ordinary quotient is finite torsion-free over
\(\Lambda\simeq\mathcal O[[X_1,\ldots,X_{nd}]]\), every nonzero
element of \(\Lambda\) is a nonzerodivisor on \(B/J\) and therefore
misses its minimal primes. Hence \(\Lambda\hookrightarrow C\) is
integral, and lying over, going up and incomparability give
\(\dim C=1+nd\). Absolute irreducibility persists at the generic
point. Choose \(n^2\) group matrices spanning the full matrix algebra
at the seed. Their trace Gram determinant is nonzero there, since
the matrix trace pairing is nondegenerate. Its value for the
universal determinant is consequently nonzero in \(C\).
Over the algebraic closure of \(\operatorname{Frac} C\), the
determinant realization and untwisting theorem gives the framed
point. It uses the condition that the determinant kills the specified
auxiliary inertia modulo \(J\). The Gram determinant shows its
image algebra has dimension \(n^2\).
The reverse map from the framed coefficient ring to \(C\) composed
with the determinant link is \(B\twoheadrightarrow C\), by uniqueness
on the determinant and ordered local characters. It is surjective,
so its kernel has quotient dimension \(1+nd\) and lies under the
framed generic point. This is the seed dimension needed for curve
selection. In the actual frame \(d=\sum_v d_v>3C_0+7\) and every
\(d_v>5\), with \(C_0\) the cohomological bound in its hypotheses.
The arithmetic transfers at separated curves are invoked in degree
two first and then in degree three, where the additional exceptional-chain
transfer is used.
\src{src/Development/AutomorphicArithmetic/AutomorphicVarieties/CurveDVRSpecialization/SeedOf.lean}{501}{The chosen minimal prime and untwisting},
\src{src/Development/AutomorphicArithmetic/AutomorphicVarieties/CurveDVRSpecialization/SeedOf.lean}{525}{The framed point and quotient dimension},
\src{src/Development/AutomorphicArithmetic/AutomorphicVarieties/FrobeniusSmoothVectors/ContinuousCoboundaries2.lean}{212}{The strict degree inequalities}
and \src{src/Development/AutomorphicArithmetic/AutomorphicVarieties/CurveDVRSpecialization/AbsMain.lean}{21}{The arithmetic transfers}.

\paragraph{The quartic seed and the four links.}
In the rank-four case-three tower, the input is a good host datum
with \(0<e<t\) in its third character case. The integral test
\((r,u)\) has the prescribed residual matrices and similitude,
the stated global ramification bound and kernel on inertia away
from \(\ell\), and a symplectic flag matrix \(u\) above \(\ell\).
Its ordered diagonal characters agree with the prescribed
characters on inertia and reduce to them on the decomposition
group. The test is defined over a small coefficient ring.
On the fraction field the case under consideration is either
irreducible or the specified nonsplit block extension with full
two-dimensional block algebra. These are the hypotheses on the
test to which the propagation is applied.
\src{src/Development/AutomorphicArithmetic/AutomorphicVarieties/WeakShimuraTheory/NearSplit.lean}{59}{Host and integral test}
and \src{src/Development/AutomorphicArithmetic/AutomorphicVarieties/WeakShimuraTheory/RedIdeal.lean}{65}{The block-extension alternative}.

Start with its actual lattice shape on the first tensor-factor tower \(T_0\).
The seed theorem supplies a rational prime \(w\) split in the
specified tower sense and a residual framed point \(s_0\), initially
allowing \(w\) in its ramification set. Its assertion at higher
levels is universal. For every allowable \(T\ge T_0\) at which
all continuous order-\(\ell\) characters of \(G_{\mathbb Q}\)
unramified away from \(w\) become trivial on \(G_{F_T}\), it
supplies an integral pair \(Q_T\) over a finite coefficient DVR.
Its reduction lies over \(s_0\), has the specified root-entry
nonvanishings, and obeys the frame with the auxiliary set now empty.
Its generic fibre is irreducible on the designated subgroup.
After restriction to the CM field of \(T\) and a specified scalar
twist it is attached to a regular algebraic cuspidal degree-four
representation. Its consecutive weight differences are \(2n_i\),
\(n_i\ge2\). The source uses a tensor product of two planes. It
chooses \(1\le j\le9\), a positive \(m\), and \(k\ge2(9m)+3\).
The three half-gaps in this construction are
\(2mj,\ k-1-2mj,\ 2mj\), all at least two. The restriction
condition at \(w\) kills the auxiliary character on \(G_{F_T}\),
which is where the integral pairing without the extra allowed prime
is checked.
\src{src/Development/AutomorphicArithmetic/AutomorphicVarieties/WeakShimuraTheory/NearSplit.lean}{175}{Condition at the auxiliary prime},
\src{src/Development/AutomorphicArithmetic/AutomorphicVarieties/ModularityArgument/ThresholdArguments/DenseThresholdCases/BirthSeedIV.lean}{34}{Tensor seed and size of the gaps},
\src{src/Development/AutomorphicArithmetic/AutomorphicVarieties/ModularityArgument/ThresholdArguments/DenseThresholdCases/BirthSeedIV.lean}{48}{Residual and integral pair on each allowed layer}
and \src{src/Development/AutomorphicArithmetic/AutomorphicVarieties/ModularityArgument/ThresholdArguments/DenseThresholdCases/SeedOf4P.lean}{50}{Weights and regular cuspidal attachment}.

For the given test lattice the cohomological construction chooses a
finite set of split auxiliary primes containing \(w\) and then
\(T_1\ge T_0\) killing all the indicated order-\(\ell\) characters
for each auxiliary prime separately. This condition
persists under further restriction of Galois groups. The cohomology
bound supplies \(C\) over the resulting tame tower, before the choice
of a deep layer \(T\ge T_1\) with \([F_T:\mathbb Q]\ge30(C+1)^2\).
The seed theorem consequently applies at this very layer. Local
conditions on the restricted test can be checked in coordinates.
For a selected \(v\mid\ell\) choose \(g_v\) carrying the two chosen
prolongations into agreement. In the transported flag the \(i\)th
character at \(s\in I_v\) is the earlier \(i\)th character at
\(g_v^{-1}s g_v\) (using restriction when viewing \(s\) over
\(\mathbb Q\)). This equals the prescribed global character
\(\psi_i(s)\), since a character is invariant under conjugation.
To see why the flag has an integral normalized basis, put
\(Y=r(g_v)n_0\) in the small local coefficient ring, with \(n_0\)
the original lower unipotent flag matrix. The prescribed residual
support gives \(\bar Y_{11}=\bar r(g_v)_{11}\),
\(\bar Y_{12}=0\) and \(\bar Y_{22}=\bar r(g_v)_{22}\) in one-based
indices. Both leading principal minors of orders one and two are
therefore units. The similitude multiplier \(\mu\) is a unit.
The third leading minor is \(\det(Y)(Y^{-1})_{44}\), a unit because
\((Y^{-1})_{44}=Y_{11}/\mu\) in the chosen antidiagonal
symplectic basis. Gaussian elimination over the local ring now
gives \(Y=NU\) with \(N\) lower unipotent and \(U\) upper triangular
invertible. Let \(X^\dagger\) denote the symplectic adjoint.
The automorphism \(X\mapsto(X^\dagger)^{-1}\) preserves lower
unipotent and upper triangular matrices and takes \(Y\) to
\(\mu^{-1}Y\). Uniqueness of triangular factorization gives
\((N^\dagger)^{-1}=N\), so the new flag basis is symplectic as
required. Restricting the multiplier identity and the
specified inertia kernels then gives the framed pair on \(T\).
\src{src/Development/AutomorphicArithmetic/AutomorphicVarieties/WeakShimuraTheory/RootsOfUnityShapes/ShapeQ.lean}{6}{The residual pivots},
\src{src/Development/AutomorphicArithmetic/AutomorphicVarieties/WeakShimuraTheory/CompactIntegralFields/RestrQ.lean}{77}{Transport and integral normalization of the flag}
and \src{src/Development/AutomorphicArithmetic/AutomorphicVarieties/ModularityArgument/ThresholdArguments/DenseThresholdCases/PromodularityCaseThreeOfLemmas.lean}{138}{Choice of lattice, cohomology bound, layer and seed}.

At that layer the explicit residual chain, in the notation of the
construction, is
\[
 q_0\ \longleftrightarrow\ q_3\ \longleftrightarrow\ q_4\
 \longleftrightarrow\ q_2\ \longleftrightarrow\ \overline Q_T .
\]
The first vertex is the restricted test, identified by uniqueness
over its chosen lattice shape. The four edges alternate between
two splittings into pairs of coordinate planes, with indices
\((1,3)\mid(2,4)\) on the first and third edges and
\((1,2)\mid(3,4)\) on the second and fourth. Here is the
reason they give actual lines in the representation scheme.
In the indicated order the two endpoint matrices are block lower
triangular with identical diagonal blocks. The same property
holds for their local flag matrices. Put
\(X(t)=X_0+t(X_1-X_0)\), with \(t\) an indeterminate. For two
such pairs, \((X_1-X_0)(Y_1-Y_0)=0\), since each difference has
only a lower left block. Expansion gives
\[
 X(t)Y(t)=X_0Y_0+t(X_1Y_1-X_0Y_0).
\]
The product pair again has the same diagonal blocks. The
symplectic adjoint is linear and preserves this block condition
in both orders. The formula consequently interpolates the group
law, the similitude identity and the conjugated flag matrices.
Equal diagonal entries, prescribed zeros and the inertia kernels
persist along the line. Evaluating at \(t=0,1\) recovers the
two framed endpoints in one irreducible image.
\src{src/Development/AutomorphicArithmetic/AutomorphicVarieties/WeakShimuraTheory/CompactIntegralFields/AffSmulOne.lean}{43}{First block decomposition}
and \src{src/Development/AutomorphicArithmetic/AutomorphicVarieties/WeakShimuraTheory/CompactIntegralFields/GlobalTheory.lean}{131}{Second block decomposition and flag matrices}.

In fact each nonconstant link is a closed affine line in the framed
special fibre. Write \(A=X.A\), \(B=X.B\) and \(k\) for the
residue field of the common coefficient DVR. The representing map
\(A\to k[t]\) kills \(\mathfrak m_B A\) by the line clause of
the representing space. If the framed endpoints differ, some
matrix or flag entry maps to \(c+dt\), \(d\ne0\). The coefficient
DVR maps onto \(k\), so the image contains \(t\). Thus the map is
surjective and its kernel lies below both point ideals. Every proper
coefficient ideal, including \(J\), vanishes scheme-theoretically
on all four links and on their formal arcs. A minimal prime of
\(A\) below the kernel supplies one ambient component through both
endpoints. This gives the transfer from \emph{all} components at
one end to \emph{some} component at the other.
\src{src/Development/AutomorphicArithmetic/AutomorphicVarieties/WeakShimuraTheory/RedIdeal.lean}{233}{The representing map for the residual line}
and \src{src/Development/AutomorphicArithmetic/AutomorphicVarieties/WeakShimuraTheory/CompactIntegralFields/AlgHomResUnique.lean}{56}{The two directions of component transfer}.

Even seed local vanishing and a finite weight fibre do not by
themselves extend this conclusion to a chosen test curve. For an
explicit calculation over a complete DVR \(O\) of uniformizer
\(\pi\) and residue field \(k\), take
\[
 B=O[[u,v]]/(uv),\qquad A=B[t]/(ut),\qquad J=(u).
\]
With weight parameter \(v\), the fibre \(B/(J,\pi,v)=k\).
The integral seed \((u,v,t)=(0,\pi,1)\) satisfies the stated
coefficient local vanishing at its generic point: every prime
below its kernel misses \(v\), hence contains \(u\). Its reduction
and \((0,0,0)\) lie on the closed line \(A/\mathfrak m_B A=k[t]\).
However the equal-characteristic test
\((u,v,t)=(z,0,0)\) in \(k[[z]]\) specializes to the latter
endpoint and pulls \(J\) back to \((z)\). The two minimal primes
of \(A\) there are \((u)\) and \((v,t)\); only the first contains
\(J\). This identifies the precise additional conclusion required
at a vertex with more than one component.

The root-entry and support tests at the five vertices verify the
good-point hypotheses of the component theorem. Enlarge the coefficient
DVR once to contain the seed coefficients and realize the residual
chain, a ring of definition of the test, and the lower coefficient
extension for the tower ideal over a common coefficient base.
The universal lifts now give points of one representing space. Local
vanishing at the generic seed puts some component at
\(\overline Q_T\) in \(V(J)\). The good-point component theorem at this
vertex gives all components there. The last line then gives some
component at \(q_2\), and the same theorem gives all components
there. Repeating backwards along the three remaining lines gives all
components at \(q_0\). Thus the specific generization representing
the test kills \(J\). This argument uses the cited arithmetic
construction of the lines, the seed local-vanishing theorem and
the good-point component theorem for the case-three hypotheses.
\src{src/Development/AutomorphicArithmetic/AutomorphicVarieties/ModularityArgument/ThresholdArguments/DenseThresholdCases/CuPChainIV.lean}{25}{Enlargement killing the auxiliary primes},
\src{src/Development/AutomorphicArithmetic/AutomorphicVarieties/ModularityArgument/ThresholdArguments/DenseThresholdCases/CuPChainIV.lean}{64}{The five good vertices and the four lines},
\src{src/Development/AutomorphicArithmetic/AutomorphicVarieties/ModularityArgument/ThresholdArguments/DenseThresholdCases/PromodularityCaseThreeOfLemmas.lean}{59}{Common coefficient extension}
and \src{src/Development/AutomorphicArithmetic/AutomorphicVarieties/ModularityArgument/ThresholdArguments/DenseThresholdCases/PromodularityCaseThreeOfLemmas.lean}{117}{Application to the chosen generization}.

The marked ternary polarization alone gives
\(\rho^c\simeq\rho^\vee\); it does not supply the invariant form
or the nontrivial projective involution in
Proposition~\ref{hr:orth}. In degree four
Proposition~\ref{hr:ind} requires a real quadratic self-twist,
and in equal characteristic it also requires the unramified
condition there. In characteristic zero Corollary~\ref{hr:imprimitive}
leaves the strongly irreducible quartic points. In equal
characteristic Proposition~\ref{hr:case3strong} also leaves those
points, together with the possible ramified real quadratic induction
in the second host case. The even two-dimensional block of
Lemma~\ref{spreal:even} and the complementary ternary cases also
remain in the general result.
